\documentclass[12pt]{article}
\usepackage[utf8]{inputenc}
\usepackage[T1]{fontenc}
\usepackage[spanish,mexico]{babel}
\usepackage{amsmath,amssymb}
\usepackage{geometry}
\usepackage{enumitem}
\usepackage{mathtools}
\geometry{letterpaper,margin=2.3cm}
\setlength{\parindent}{0pt}
\setlist[enumerate]{itemsep=.55em,topsep=.4em}

\begin{document}

\section*{Ejercicios 1.1}

\textit{En los ejercicios 1 a 4, determine si el conjunto es una función. Si es una función determine su dominio.}

\begin{enumerate}
\item
\begin{enumerate}[label=(\alph*)]
\item $\{(x,y)\mid y=\sqrt{x-4}\}$
\item $\{(x,y)\mid y=\sqrt{x^2-4}\}$
\item $\{(x,y)\mid y=\sqrt{4-x^2}\}$
\item $\{(x,y)\mid x^2+y^2=4\}$
\end{enumerate}

\item
\begin{enumerate}[label=(\alph*)]
\item $\{(x,y)\mid y=\sqrt{x+1}\}$
\item $\{(x,y)\mid y=\sqrt{x^2-1}\}$
\item $\{(x,y)\mid y=\sqrt{1-x^2}\}$
\item $\{(x,y)\mid x^2+y^2=1\}$
\end{enumerate}

\item
\begin{enumerate}[label=(\alph*)]
\item $\{(x,y)\mid y=x^2\}$
\item $\{(x,y)\mid x=y^2\}$
\item $\{(x,y)\mid y=x^3\}$
\item $\{(x,y)\mid x=y^3\}$
\end{enumerate}

\item
\begin{enumerate}[label=(\alph*)]
\item $\{(x,y)\mid y=(x-1)^2+2\}$
\item $\{(x,y)\mid x=(y-2)^2+1\}$
\item $\{(x,y)\mid y=(x+2)^3-1\}$
\item $\{(x,y)\mid x=(y+1)^3-2\}$
\end{enumerate}

\item Dada $f(x)=2x-1$, determine
\begin{enumerate}[label=(\alph*)]
\item $f(3)$; \item $f(-2)$; \item $f(0)$; \item $f(a+1)$; \item $f(x+1)$;
\item $f(2x)$; \item $2f(x)$; \item $f(x+h)$; \item $f(x)+f(h)$;
\item $\dfrac{f(x+h)-f(x)}{h}$, $h\ne0$.
\end{enumerate}

\item Dada $f(x)=\dfrac{3}{x}$, calcule
\begin{enumerate}[label=(\alph*)]
\item $f(1)$; \item $f(-3)$; \item $f(6)$; \item $f\!\left(\dfrac13\right)$;
\item $f\!\left(\dfrac3a\right)$; \item $f\!\left(\dfrac3x\right)$;
\item $\dfrac{f(3)}{f(x)}$; \item $f(x-3)$; \item $f(x)-f(3)$;
\item $\dfrac{f(x+h)-f(x)}{h}$, $h\ne0$.
\end{enumerate}

\item Dada $f(x)=2x^2+5x-3$, determine
\begin{enumerate}[label=(\alph*)]
\item $f(-2)$; \item $f(-1)$; \item $f(0)$; \item $f(3)$; \item $f(h+1)$;
\item $f(2x^2)$; \item $f(x^2-3)$; \item $f(x+h)$; \item $f(x)+f(h)$;
\item $\dfrac{f(x+h)-f(x)}{h}$, $h\ne0$.
\end{enumerate}

\item Dada $g(x)=3x^2-4$, calcule
\begin{enumerate}[label=(\alph*)]
\item $g(-4)$; \item $g\!\left(\dfrac12\right)$; \item $g(x^2)$; \item $g(3x^2-4)$;
\item $g(x-h)$; \item $g(x)-g(h)$;
\item $\dfrac{g(x+h)-g(x)}{h}$, $h\ne0$.
\end{enumerate}

\item Dada $F(x)=\sqrt{x+9}$, encuentre
\begin{enumerate}[label=(\alph*)]
\item $F(x+9)$; \item $F(x^2-9)$; \item $F(x^4-9)$; \item $F(x^2+6x)$;
\item $F(x^4-6x^2)$; \item $\dfrac{F(x+h)-F(x)}{h}$, $h\ne0$.
\end{enumerate}

\item Dada $G(x)=\sqrt{4-x}$, determine
\begin{enumerate}[label=(\alph*)]
\item $G(4-x)$; \item $G(4-x^2)$; \item $G(4-x^4)$; \item $G(4x-x^2)$;
\item $G(-x^4-4x^2)$; \item $\dfrac{G(x+h)-G(x)}{h}$, $h\ne0$.
\end{enumerate}
\end{enumerate}

\textit{En los ejercicios 11 a 46, dibuje a mano la gráfica de la función y determine su dominio y su contradominio.}

\begin{enumerate}[resume]
\item $f(x)=3x-1$
\item $g(x)=4-x$
\item $F(x)=2x^2$
\item $G(x)=x^2+2$
\item $g(x)=5-x^2$
\item $f(x)=(x-1)^2$
\item $G(x)=\sqrt{x-1}$
\item $F(x)=\sqrt{9-x}$
\item $f(x)=\sqrt{x^2-4}$
\item $g(x)=\sqrt{4-x^2}$
\item $g(x)=\sqrt{9-x^2}$
\item $f(x)=\sqrt{x^2-1}$
\item $h(x)=|x-3|$
\item $H(x)=|5-x|$
\item $F(x)=|3x+2|$
\item $G(x)=\dfrac{x^2-4}{x-2}$
\item $H(x)=\dfrac{x^2-25}{x+5}$
\item $f(x)=\dfrac{2x^2+7x+3}{x+3}$
\item $f(x)=\dfrac{x^2-4x+3}{x-1}$
\item $g(x)=\dfrac{(x^2-4)(x-3)}{x^2-x-6}$
\item $f(x)=\begin{cases}-2,&x\le3,\\2,&3<x.\end{cases}$
\item $g(x)=\begin{cases}-4,&x<-2,\\-1,&-2\le x\le2,\\3,&2<x.\end{cases}$
\item $g(x)=\begin{cases}2x-1,&x\ne2,\\0,&x=2.\end{cases}$
\item $f(x)=\begin{cases}3x+2,&x\ne1,\\8,&x=1.\end{cases}$
\item $F(x)=\begin{cases}x^2-4,&x\ne3,\\-2,&x=3.\end{cases}$
\item $G(x)=\begin{cases}9-x^2,&x\ne-3,\\4,&x=-3.\end{cases}$
\item $G(x)=\begin{cases}1-x^2,&x<0,\\3x+1,&0\le x.\end{cases}$
\item $F(x)=\begin{cases}x^2-4,&x<3,\\2x-1,&3\le x.\end{cases}$
\item $g(x)=\begin{cases}6x+7,&x\le-2,\\4-x,&-2<x.\end{cases}$
\item $f(x)=\begin{cases}x-2,&x\le0,\\x^2+1,&0<x.\end{cases}$
\item $h(x)=\begin{cases}x+3,&x<-5,\\\sqrt{25-x^2},&-5\le x\le5,\\3-x,&5<x.\end{cases}$
\item $H(x)=\begin{cases}\dfrac{x+2}{\sqrt{16-x^2}},&x\le-4,\\2-x,&-4<x<4,\\2-x,&4\le x.\end{cases}$
\item $F(x)=\dfrac{x^3-2x^2}{x-2}$
\item $G(x)=\dfrac{x^3+3x^2}{x+3}$
\item $f(x)=\lfloor x-4\rfloor$
\item $g(x)=\lfloor x+2\rfloor$
\end{enumerate}

\begin{enumerate}[start=47]
\item
\begin{enumerate}[label=(\alph*)]
\item Dibuje la gráfica de la función escalón (o salto) unitario denotada por $U$ y definida por
\[
U(x)=\begin{cases}0,&x<0,\\1,&0\le x.\end{cases}
\]
Defina cada una de las siguientes funciones a trozos y dibuje sus gráficas:
\item $U(x-1)$; \item $U(x)-1$; \item $U(x)-U(x-1)$.
\end{enumerate}

\item Defina cada una de las siguientes funciones a trozos y dibuje sus gráficas, donde $U$ es la función escalón unitario definida en el ejercicio 47:
\begin{enumerate}[label=(\alph*)]
\item $x\,U(x)$; \item $(x+1)U(x+1)$; \item $(x+1)U(x+1)-xU(x)$.
\end{enumerate}

\item
\begin{enumerate}[label=(\alph*)]
\item Dibuje la gráfica de la función signo denotada por $\operatorname{sgn}$ y definida por
\[
\operatorname{sgn}x=\begin{cases}-1,&x<0,\\0,&x=0,\\1,&0<x.\end{cases}
\]
$\operatorname{sgn}(x)$ se lee ``signo de $x$''. Defina cada una de las siguientes funciones a trozos y dibuje sus gráficas:
\item $x\,\operatorname{sgn}(x)$; \item $2-x\,\operatorname{sgn}(x)$; \item $x-2\,\operatorname{sgn}(x)$.
\end{enumerate}

\item Defina cada una de las siguientes funciones a trozos, donde $\operatorname{sgn}$ es la función signo definida en el ejercicio 49:
\begin{enumerate}[label=(\alph*)]
\item $\operatorname{sgn}(x+1)$; \item $\operatorname{sgn}(x-1)$; \item $\operatorname{sgn}(x+1)-\operatorname{sgn}(x-1)$.
\end{enumerate}

\item La gráfica de la función $f$ de la figura se parece a la letra W. Defina $f(x)$ a trozos.
\item La gráfica de la función $f$ de la figura se parece a la letra M. Defina $f(x)$ a trozos.
\item En la figura, la gráfica se parece a la letra X y es la gráfica de dos funciones $f_1$ y $f_2$ trazadas en el rectángulo de inspección de $[-1,1]$ por $[-1,1]$. Defina $f_1(x)$ y $f_2(x)$.
\item Existen tres funciones $f_1$, $f_2$ y $f_3$ cuyas gráficas trazadas simultáneamente en el rectángulo de inspección de $[-1,1]$ por $[-1,1]$ se parecen a la letra Z. Defina $f_1(x)$, $f_2(x)$ y $f_3(x)$.
\end{enumerate}

\textit{En los ejercicios 55 a 58, haga lo siguiente: (a) defina la función a trozos sin emplear las barras de valor absoluto; (b) dibuje la gráfica de la función definida en el inciso (a); (c) apoye sus respuestas a los incisos (a) y (b) trazando la gráfica de la función.}

\begin{enumerate}[start=55]
\item $f(x)=|x^2-1|$
\item $g(x)=|4-x^2|$
\item $g(x)=|x|\,|5-x|$
\item $f(x)=|x|\,|x-3|$
\end{enumerate}

\newpage
\section*{Ejercicios 1.2}

\textit{En los ejercicios 1 a 10, defina las siguientes funciones y determine el dominio de la función resultante: (a) $f+g$; (b) $f-g$; (c) $f\cdot g$; (d) $f/g$; (e) $g/f$.}

\begin{enumerate}
\item $f(x)=x-5$; $g(x)=x^2-1$
\item $f(x)=\sqrt{x}$; $g(x)=x^2+1$
\item $f(x)=\dfrac{x+1}{x-1}$; $g(x)=\dfrac1x$
\item $f(x)=\sqrt{x}$; $g(x)=4-x^2$
\item $f(x)=\sqrt{x}$; $g(x)=x^2-1$
\item $f(x)=|x|$; $g(x)=|x-3|$
\item $f(x)=x^2+1$; $g(x)=3x-2$
\item $f(x)=\sqrt{x-4}$; $g(x)=x^2-4$
\item $f(x)=\dfrac1{x+1}$; $g(x)=\dfrac{x}{x-2}$
\item $f(x)=x^2$; $g(x)=\dfrac1{\sqrt{x}}$
\end{enumerate}

\textit{En los ejercicios 11 a 14, para las funciones $f$ y $g$ y el número $c$, obtenga $(f\circ g)(c)$ mediante dos métodos: (a) calcule $g(c)$ y utilice este número para determinar $f(g(c))$; (b) determine $(f\circ g)(x)$ y emplee este valor para calcular $(f\circ g)(c)$.}

\begin{enumerate}[start=11]
\item $f(x)=3x^2-4x$; $g(x)=2x-5$; $c=4$
\item $f(x)=\sqrt{x^2-36}$; $g(x)=x^2-3x$; $c=5$
\item $f(x)=\dfrac1{x-1}$; $g(x)=\dfrac2{x^2+1}$; $c=\dfrac12$
\item $f(x)=\dfrac{2\sqrt{x}+3}{x}$; $g(x)=\dfrac{2x+5}{x^4}$; $c=-2$
\end{enumerate}

\textit{En los ejercicios 15 a 24, defina las siguientes funciones y determine el dominio de la función compuesta: (a) $f\circ g$; (b) $g\circ f$; (c) $f\circ f$; (d) $g\circ g$.}

\begin{enumerate}[start=15]
\item $f(x)=x-2$; $g(x)=x+7$
\item $f(x)=3-2x$; $g(x)=6-3x$
\item Las funciones del ejercicio 1.
\item Las funciones del ejercicio 2.
\item $f(x)=\sqrt{x+2}$; $g(x)=x^2-2$
\item $f(x)=x^2-1$; $g(x)=\dfrac1x$
\item $f(x)=\dfrac1x$; $g(x)=\sqrt{x}$
\item $f(x)=\sqrt{x}$; $g(x)=-\dfrac1x$
\item $f(x)=|x|$; $g(x)=|x+2|$
\item $f(x)=\sqrt{x^2-1}$; $g(x)=\sqrt{x-1}$
\end{enumerate}

\textit{En los ejercicios 25 y 26 defina las siguientes funciones y determine el dominio de las funciones resultantes: (a) $f(x^2)$; (b) $[f(x)]^2$; (c) $f\circ f(x)$; (d) $f\circ f(-x)$.}

\begin{enumerate}[start=25]
\item $f(x)=\sqrt{x}$
\item $f(x)=\dfrac1{x-1}$
\end{enumerate}

\textit{En los ejercicios 27 a 32, exprese $h$ como composición de las dos funciones $f$ y $g$ en dos formas.}

\begin{enumerate}[start=27]
\item $h(x)=\sqrt{x^2-4}$
\item $h(x)=(9+x^2)^{-2}$
\item $h(x)=\left(\dfrac1{x-2}\right)^3$
\item $h(x)=\dfrac4{\sqrt[3]{x^3+3}}$
\item $h(x)=(x^2+4x-5)^4$
\item $h(x)=\sqrt{|x|+4}$
\end{enumerate}

\textit{En los ejercicios 33 a 38, trace en la graficadora la gráfica de la función y a partir de la gráfica conjeture si la función es par, impar o de ninguno de estos dos tipos. Después confirme su conjetura analíticamente.}

\begin{enumerate}[start=33]
\item \begin{enumerate}[label=(\alph*)]\item $f(x)=2x^4-3x^2+1$ \item $g(x)=5x^5+1$\end{enumerate}
\item \begin{enumerate}[label=(\alph*)]\item $f(x)=x^2+2x+2$ \item $g(x)=x^6-1$\end{enumerate}
\item \begin{enumerate}[label=(\alph*)]\item $f(x)=5x^3-7x$ \item $g(x)=|x|$\end{enumerate}
\item \begin{enumerate}[label=(\alph*)]\item $f(x)=4x^5+3x^3$ \item $g(x)=x^3+1$\end{enumerate}
\item \begin{enumerate}[label=(\alph*)]\item $f(x)=\sqrt[3]{x}$ \item $g(x)=5x^4-4$\end{enumerate}
\item \begin{enumerate}[label=(\alph*)]\item $f(x)=\dfrac{|x|}{x}$ \item $g(x)=2|x|+3$\end{enumerate}
\end{enumerate}

\textit{En los ejercicios 39 y 40, determine analíticamente si la función es par, impar o de ninguno de estos dos tipos.}

\begin{enumerate}[start=39]
\item \begin{enumerate}[label=(\alph*)]\item $f(y)=\dfrac{y^3-y}{y^2+1}$ \item $g(r)=\dfrac{r^2-1}{r^2+1}$ \item $f(x)=\dfrac{|x|}{x^2+1}$\end{enumerate}
\item \begin{enumerate}[label=(\alph*)]\item $h(x)=\dfrac{x^2-5}{2x^3+x}$ \item $g(z)=\dfrac{z-1}{z+1}$ \item $f(x)=\begin{cases}-1,&x<0,\\1,&0\le x.\end{cases}$\end{enumerate}
\end{enumerate}

\textit{En los ejercicios 41 a 44, haga lo siguiente: (a) defina $f(x)$, sin las barras de valor absoluto, en los intervalos indicados. (b) Apoye la respuesta del inciso (a), gráficamente trazando la gráfica de $f$ en la graficadora a partir de la ecuación dada. (c) A partir de la gráfica del inciso (b), establezca si la función es par, impar o de ninguno de estos dos tipos. (d) Confirme la respuesta del inciso (c) analíticamente a partir de la ecuación dada.}

\begin{enumerate}[start=41]
\item $f(x)=\dfrac{|x|}{x}$; $(-\infty,0)$, $(0,+\infty)$
\item $f(x)=x|x|$; $(-\infty,0)$, $[0,+\infty)$
\item $f(x)=|x-2|-|x+2|$; $(-\infty,-2)$, $[-2,2)$, $[2,+\infty)$
\item $f(x)=\dfrac{|x+1|-|x-1|}{x}$; $(-\infty,-1)$, $[-1,0)$, $(0,1]$, $(1,+\infty)$
\end{enumerate}

\begin{enumerate}[start=45]
\item ¿Es conmutativa la composición de dos funciones? Es decir, si $f$ y $g$ son dos funciones cualesquiera, ¿son iguales $(f\circ g)(x)$ y $(g\circ f)(x)$? Justifique su respuesta proporcionando un ejemplo.
\end{enumerate}

Si $f$ y $g$ son dos funciones tales que $(f\circ g)(x)=x$ y $(g\circ f)(x)=x$, entonces se dice que $f$ y $g$ son inversas una de la otra. En los ejercicios 46 a 50, demuestre que $f$ y $g$ son inversas una de la otra.

\begin{enumerate}[start=46]
\item $f(x)=2x-3$ y $g(x)=\dfrac{x+3}{2}$
\item $f(x)=\dfrac1{x+1}$ y $g(x)=\dfrac{1-x}{x}$
\item $f(x)=x^2$, $x\ge0$, y $g(x)=\sqrt{x}$
\item $f(x)=x^2$, $x\le0$, y $g(x)=-\sqrt{x}$
\item $f(x)=(x-1)^3$ y $g(x)=1+\sqrt[3]{x}$
\item La función escalón unitario $U$ y la función signo $\operatorname{sgn}$ se definieron en los ejercicios 47 y 49, respectivamente, de la sección 1.1. (a) Defina $\operatorname{sgn}(U(x))$ y dibuje la gráfica. (b) Defina $U(\operatorname{sgn}(x))$ y dibuje la gráfica.
\item Demuestre que si $f$ y $g$ son funciones impares, entonces $(f+g)$ y $(f-g)$ también son funciones impares, mientras que $fg$ y $f/g$ son funciones pares.
\item Determine si la función compuesta $f\circ g$ es par o impar en cada uno de los casos siguientes: (a) $f$ y $g$ son impares; (b) $f$ es par y $g$ es impar; (c) $g$ es par.
\item Encuentre fórmulas para $(f\circ g)(x)$ si
\[
f(x)=\begin{cases}0,&x<0,\\2x,&0\le x\le1,\\0,&1<x,\end{cases}
\qquad
y\qquad
g(x)=\begin{cases}1,&x<0,\\\frac12x,&0\le x\le1,\\1,&1<x.\end{cases}
\]
Dibuje las gráficas de $f$, $g$ y $f\circ g$.
\item Encuentre fórmulas para $(g\circ f)(x)$ a partir de las funciones del ejercicio 54. Dibuje la gráfica de $g\circ f$.
\item Si $f(x)=x^2+2x+2$, encuentre dos funciones $g$ para las cuales $(f\circ g)(x)=x^2-4x+5$.
\item Si $f(x)=x^2$, encuentre dos funciones $g$ para las cuales $(f\circ g)(x)=4x^2-12x+9$.
\item Demuestre que si $f$ y $g$ son dos funciones lineales, entonces $f\circ g$ es una función lineal.
\item Existe una función cuyo dominio es el conjunto de todos los números reales que es a la vez par e impar. ¿Cuál es esa función? Demuestre que es única esta función.
\item Suponga que $f(x)=\dfrac1x$, $g(x)=-\dfrac1x$ y $h(x)=-x$. Demuestre que $(f\circ g)(x)=(g\circ f)(x)$ y explique por qué $f\circ g$ ni $g\circ f$ son la misma que $h$.
\item Trace en la graficadora las gráficas de las dos funciones $F$ y $G$ definidas por
\[
F(x)=\frac{\sqrt{x+1}}{\sqrt{x-4}}
\qquad\text{y}\qquad
G(x)=\sqrt{\frac{x+1}{x-4}}.
\]
Observe que $F$ es la misma función que $f/g$ del ejemplo 1(d). Explique por qué las gráficas de $F$ y $G$ no son las mismas y, consecuentemente, las funciones no son iguales.
\end{enumerate}

\end{document}
